\section{Discussion}
\label{sec:discussion}

\paragraph{What the margin view buys.}
Writing rejection as a margin did not produce a method that beats GCE on symmetric noise. It did
make three properties of selection rules visible before any training run, and each was then
measured. First, \emph{relative} statistics carry a hidden rate. A z-score, a rank or a fixed
budget rejects the same share of a clean batch as of a corrupted one (\cref{prop:blind}). That
is why the batch-relative margin wins exactly at one noise level and why small-loss needs the
oracle rate. Second, an \emph{absolute} statistic turns the rejection threshold into a constant of
the loss, a band of label gaps (\cref{cor:band}). The rejected share then follows the data,
including structured noise, where rate-based selection has no reason to pick the right samples.
Third, any margin family with a floor has an \emph{escape region} (\cref{prop:reject}). It
explains the threshold-shaped strength sweep and the failure of the wider-temperature variant.
It also identifies a defect of our recipe: the most confidently contradicted labels receive a
near-maximal gradient toward the wrong class. The remedy is $\beta \ge 1 + \dbase$, which makes
the band reach $g = 1$; with $\beta = 1.15$, $\tau = 0.1$ it becomes $[0.032, 1.0)$. We did not
run it, because the pre-registered protocol froze $\beta$ before Phase~C$'$, and we leave it as a
stated prediction.

\paragraph{When to reject and when to down-weight.}
The absolute rule discards every sample the model currently disagrees with. Under symmetric noise
many of those are hard but correct, and GCE, which only down-weights them, is 0.9--1.3 points
better. Under pair-flip noise, disagreement is the better signal. A fixed wrong class is learned
as a pattern, so mislabeled samples do not stand out by loss rank (small-loss memorizes half of
them). They appear to stand out by disagreement before the pattern is learned (the absolute rule
memorizes a quarter).
The same contrast appears without model selection. The heavier or more structured the noise, the
larger the absolute rule's last-epoch lead over GCE: $+5.1$ at 60\% symmetric and $+7.1$ at 40\%
pair-flip.

\paragraph{Limitations.}
One dataset, one model, one modality, and synthetic noise only. The real-noise benchmark was cut
by our own pre-registered rule, so the claims above are about 20~Newsgroups with injected noise.
Five seeds resolve differences of about one point at best: several comparisons we describe
($+4.9$ over GCE under pair-flip, $-0.95$ at 40\% symmetric) are not established at $p < 0.05$.
Small-loss is the single-network rule, not Co-teaching's two networks, and DivideMix-style
semi-supervised pipelines are out of scope, since they are training procedures rather than losses.
The baselines' tuning grid ran only at 40\% symmetric noise. Training was fixed at 8 epochs with a
constant learning rate. The suspicion margins inherit AS-Softmax's base margin, which we showed
memorizes faster than CE on its own.

\section{Conclusion}
\label{sec:conclusion}

In masked softmax, a sample's loss vanishes exactly when its label gap reaches the negative of its
margin. A negative margin is therefore sample rejection, and the choice of selection rule becomes
the choice of a margin function. The view makes falsifiable predictions, and they held in a
pre-registered study. Relative statistics are noise-blind. An absolute ``another class beats the
label'' statistic rejects a fixed, rate-free band of gaps, memorizes the fewest corrupted labels and
is the most robust rule under structured noise. A margin floor leaves an escape region that
shapes how the rule responds to its strength. The absolute rule did not clear our bar against GCE
on clean and symmetric noise, and we report it as a unifying account with measured consequences,
not as a new state of the art.
